%% kapa: exam=1 round=22 year=2011 period=2 form=A subject=accounting subject_ko=회계학 seq=35 number=35
%% source: Q-Net 감정평가사 기출문제 (cst00309) artlSeq=5036790 file=extracted/1008/5036790/2011년 제22회 감정평가사 1차시험 기출문제/2교시.hwp
%% conversion: kordoc 4.24.51 -> tools/gen1.py (auto); review=proofread-sonnet; flags=none
\begin{question}{35}
(주)서울은 20×1년에 제품A를 연간 1,500단위 생산하여 단위당 ￦400에 판매하였다. 제품A의 최대생산량은 2,000단위이며 단위당 원가는 다음과 같다.

{\small\centering
\begin{tabular}{|lr|}
\hline
ㆍ직접재료원가 & ￦120 \\
ㆍ직접노무원가 & 80 \\
ㆍ변동제조간접원가 & 20 \\
ㆍ변동판매관리비 & 30 \\
ㆍ고정판매관리비 & 20 \\
ㆍ고정제조간접원가 & 30 \\
\hline
\end{tabular}\par}

20×2년 초에 회사는 (주)한국으로부터 제품A 800단위를 단위당 ￦300에 구입하겠다는 특별주문을 받았다. (주)서울이 동 주문을 수락하면 단위당 변동판매관리비 중 ￦20이 발생하지 않으며, 기존시장에서의 판매량 300단위를 포기해야 한다. (주)서울이 특별주문 수량을 모두 수락할 경우 이익은 얼마나 증가하는가? (단, 재고는 없으며, 20×2년 원가구조는 20×1년과 동일함) \par
\begin{choices}
  \item ￦10,200
  \item ￦10,400
  \item ￦10,600
  \item ￦10,800
  \item ￦11,000
\end{choices}
\end{question}
\answer{5}
